


\section{Evaluation}
\label{sec:eval}

%The primary goal of developer-written/automated-generated test suites is to robustly and reliably detect faults should they be inadvertently introduced to a system. To do this, a test 

Unit test suites strive to comprehensively validate that the code under test behaves as expected. 
Structural metrics such as line coverage and mutation kill score are widely used to assess this, but they do not capture whether all expected behaviours are actually validated. 
In this section, we investigate whether gaps exist between the behaviours developers document and the behaviours their test suites validate.
%
%Specifically, our evaluation seeks to answer the four research questions posed in the introduction.

Specifically, our evaluation seeks to answer the following research questions:
% \rth{do we need these questions again? i like the additional context, but they're in the intro too now :). or maybe being clear is better. we can decide about this last based on space} 

\begin{itemize}
    \item \textbf{RQ1 (Feasibility): With what precision can expected behaviours be extracted from source code and documentation?} 
     Before studying behavioural gaps, we must establish whether automatically-extracted behaviours are meaningful and grounded in the documented intent of each method. This RQ evaluates the quality of behaviour extraction as a precondition for the remaining research questions.
    
    \item \textbf{RQ2 (Prevalence): Do developer-written test suites exhibit behavioural gaps in practice?} 
     This RQ investigates whether behavioural gaps exist in practice, how large they are, and what kinds of behaviours developers commonly overlook in their test suites.
    
    \item \textbf{RQ3 (Generalization): To what extent are auto\-matically-generated test suites prone to behavioural gaps?} 
     This RQ determines whether automatically-generated suites also exhibit behavioural gaps to determine whether these gaps are inherent to manual test writing practice.
    
    \item \textbf{RQ4 (Implication): Do behavioural gaps alias traditional structural metrics?} 
     This RQ examines whether behavioural gaps are captured by existing structural metrics such as line coverage and mutation kill score, or whether they represent a complementary dimension of test suite adequacy.
\end{itemize}

Experimental design decisions all have benefits and drawbacks. Following best-practice advice~\cite{design-tradeoffs}, we report the experimental design decisions and their implications along with our experimental methodology rather than a separate threats to validity section. 
The full replication package, including all identified untested behaviours is publicly available.\footnote{\toolname{}: \url{https://anonymous.4open.science/r/behavioural-coverage}}

\smallskip
\noindent\textbf{LLM Configuration.} All experiments were conducted using the Qwen3-14B model~\cite{qwen3technicalreport}, executed locally on a single NVIDIA A100 80GB GPU with recommended settings (Temperature=0.6, TopP=0.95, TopK=20, and MinP=0).\footnote{\url{https://huggingface.co/Qwen/Qwen3-14B\#best-practices}}

\smallskip
\noindent\textbf{Project Selection}
%To evaluate our the nature of behavioural gaps, we selected a set of popular Java projects from 
Evaluated projects were selected from 
Defects4J~\cite{10.1145/2610384.2628055}, a widely used benchmark in software testing research. 
We applied a systematic filtering process to ensure technical consistency and reproducibility. 
We retained only projects that: (1) contain a \texttt{pom.xml} file in the root directory, ensuring Maven-based builds; (2) are not part of the libraries used in our \toolname{} infrastructure; (3) are hosted as monorepos, meaning each repository contains only the library's code rather than a collection of different libraries; and (4) use JUnit version 4 or higher. Applying these criteria yielded ten projects. 
%As our goal is to study behavioural gaps, we restrict our analysis to 
We analyzed all public and protected documented methods that were executed by at least one test case.
A summary of the selected projects is shown in Table~\ref{tab:repositories}.

% this table contains the actual number of methods for all repo
% \begin{table*}[ht]
% \centering
% \caption{Dataset used from \texttt{Defect4J} along with their structural properties, test coverage, and mutation kill score.} 
% \label{tab:repositories}
% \begin{tabular}{lrrrrr}
%         \toprule
%         Project (Version) & \# Source Methods & \# Test Methods & SLOC & Line Cov. & Mutation Kill Score \\
%         \midrule
%         commons-cli (1.10.0-SNAPSHOT) & 551 & 461 & 1965 & 98.22\% & 93.00\% \\
%         commons-codec (1.21.1-SNAPSHOT) & 1145 & 1196 & 4972 & 94.87\% & 92.00\% \\
%         commons-compress (1.29.0-SNAPSHOT) & 4751 & 2254 & 25308 & 88.01\% & 82.00\% \\
%         commons-jxpath (1.4.1-SNAPSHOT) & 1615 & 367 & 9228 & 77.67\% & 75.00\% \\
%         commons-collections (4.5.1-SNAPSHOT) & 4824 & 1901 & 14751 & 85.89\% & 87.00\% \\
%         commons-lang (3.18.0-SNAPSHOT) & 4580 & 4483 & 16133 & 96.50\% & 88.00\% \\
%         commons-csv (1.14.2-SNAPSHOT) & 286 & 511 & 1225 & 99.59\% & 95.00\% \\
%         gson (2.13.3-SNAPSHOT) & 1039 & 1395 & 5253 & 92.21\% & 91.00\% \\
%         jfreechart (2.0.0-SNAPSHOT) & 8179 & 2299 & 51346 & 56.50\% & 59.00\% \\
%         jsoup (1.21.1) & 2141 & 1442 & 9778 & 89.62\% & 79.00\% \\
%         \bottomrule
%     \end{tabular}
% \end{table*}


Coverage was collected with JaCoCo~\cite{jacoco2026}; mutation kill scores were collected using PIT~\cite{pitest2026}. 
Most projects have high coverage and kill scores. Eight of ten projects exceed 80\% line coverage; seven projects exceed 80\% mutation kill score.
% The high kill scores (seven projects $>$ 80\%) indicate that these test suites already contain strong assertions, making them suitable for studying behavioural gaps that persist even when traditional metrics suggest high test quality.

\smallskip
\noindent\textbf{Implication:} 
These projects enable studying whether behavioural gaps exist in popular, long-lived projects that have invested significant testing effort.
%a best case for testing effort as they are popular, well-maintained libraries with high structural coverage. 
By investigating behavioural gaps in comprehensively-tested projects, we hope our findings represent a lower bound for missed behaviour.
%Behavioural gaps identified in this setting, therefore, likely represent a lower bound 
%on their prevalence in real-world projects where documentation and 
%testing investment may be lower. 
%as projects that have tested less stringently are more likely
The project selection methodology limits our analysis to one broad class of project (libraries) and one 
programming language (Java). 
Evaluating the approach on other languages and domains remains for future work.

% These projects enable our study to uncover whether behavioural gaps exist in popular, long-lived projects.
% The project's strong coverage and mutation scores indicate developers have investigated meaningful effort validating each project's expected behaviours.
% Negatively, the project selection methodology does limit our analysis to one broad class of project (libraries) and one programming language (Java).
% Evaluating the approach on other languages and domains remains for future studies. 